% Clase 11 --- Cuándo se puede descomponer y cuándo no (§6).
% Los tres casos que aparecieron en la clase, en el mismo orden: lado a lado
% (paralelo), sándwich (serie) y el contraejemplo con el que el docente cerró
% --- dos dieléctricos separando directamente +Q de -Q, donde la interacción
% entre las partes ya no es despreciable y la receta no aplica.
\begin{center}
\begin{tikzpicture}[scale=0.9]

  % =============== (a) lado a lado -> paralelo ===============
  \begin{scope}
    \draw[figred, line width=1.5pt] (-1.6,1.15) -- (1.6,1.15);
    \draw[figblue, line width=1.5pt] (-1.6,-1.15) -- (1.6,-1.15);
    \fill[figblue!12] (-1.5,-1.02) rectangle (0,1.02);
    \fill[figamber!22] (0,-1.02) rectangle (1.5,1.02);
    \draw[figgray, line width=0.6pt] (-1.5,-1.02) rectangle (1.5,1.02);
    \draw[figgray, line width=0.6pt] (0,-1.02) -- (0,1.02);
    \node[figlbl] at (-0.75,0) {$K_1$};
    \node[figlbl] at (0.75,0) {$K_2$};

    \node[figlblaux, anchor=south, align=center] at (0,1.3)
      {las dos mitades comparten\\[-0.2em] la \emph{misma} placa};

    \node[figlbl, figred, anchor=north, align=center] at (0,-1.75)
      {\textbf{paralelo}: mismo $\Delta V$};
    \node[figlbl, anchor=north] at (0,-2.55) {(a) lado a lado};
  \end{scope}

  % =============== (b) sándwich -> serie ===============
  \begin{scope}[xshift=5.6cm]
    \draw[figred, line width=1.5pt] (-1.6,1.15) -- (1.6,1.15);
    \draw[figblue, line width=1.5pt] (-1.6,-1.15) -- (1.6,-1.15);
    \fill[figblue!12] (-1.5,0) rectangle (1.5,1.02);
    \fill[figamber!22] (-1.5,-1.02) rectangle (1.5,0);
    \draw[figgray, line width=0.6pt] (-1.5,-1.02) rectangle (1.5,1.02);
    \node[figlbl] at (0,0.5) {$K_1$};
    \node[figlbl] at (0,-0.5) {$K_2$};
    % placa imaginaria en la interfaz
    \draw[figgray, dashed, line width=1pt] (-1.5,0) -- (1.5,0);
    \node[figlblaux, anchor=west] at (1.6,0) {placa imaginaria};

    \node[figlbl, figred, anchor=north, align=center] at (0,-1.75)
      {\textbf{serie}: misma $Q$};
    \node[figlbl, anchor=north] at (0,-2.55) {(b) sándwich};
  \end{scope}

  % =============== (c) el caso que no se descompone ===============
  \begin{scope}[xshift=12.6cm]
    \fill[figblue!12] (-1.5,-1.15) rectangle (0,1.15);
    \fill[figamber!22] (0,-1.15) rectangle (1.5,1.15);
    \draw[figgray, line width=0.6pt] (-1.5,-1.15) rectangle (1.5,1.15);
    \draw[figgray, line width=0.6pt] (0,-1.15) -- (0,1.15);
    \node[figlbl] at (-0.9,0.75) {$K_1$};
    \node[figlbl] at (0.9,0.75) {$K_2$};

    \node[figpos, minimum size=8pt] at (-0.75,-0.35) {};
    \node[figlbl, figred, anchor=north] at (-0.75,-0.5) {$+Q$};
    \node[figneg, minimum size=8pt] at (0.75,-0.35) {};
    \node[figlbl, figblue, anchor=north] at (0.75,-0.5) {$-Q$};
    % líneas de campo que cruzan la frontera
    \draw[figvecres, line width=0.8pt] (-0.6,-0.2) .. controls (-0.3,0.35) and (0.3,0.35) .. (0.6,-0.2);

    \node[figlbl, figred, anchor=north, align=center] at (0,-1.75)
      {\textbf{no} se separa};
    \node[figlbl, anchor=north, align=center] at (0,-2.55)
      {(c) cargas enfrentadas\\[-0.2em] a través de la frontera};
  \end{scope}

\end{tikzpicture}\\[0.5em]
{\small\itshape La descomposición de (a) y (b) no es un truco de notación: se
apoya en que los efectos de borde son despreciables, incluso en la frontera
entre los dos dieléctricos. Bajo esa hipótesis las dos mitades \emph{no
interactúan}, y da igual que estén pegadas o separadas ---por eso (a) es
literalmente un paralelo, con las dos regiones entre las mismas placas
conductoras y por lo tanto al mismo potencial---. Cada región queda con un solo
dieléctrico, que es la condición para poder usar $\Phi_E = Q_{\text{libre}}/
\varepsilon$; recién después se combinan las capacidades. En (c) esa hipótesis
se cae: las líneas de campo cruzan la frontera y la interacción entre las dos
regiones es del mismo orden que el problema de cada una por separado, así que no
hay dos subproblemas independientes que resolver. Ese caso se trata en el curso
de Electromagnetismo.}
\end{center}
